\documentclass[12pt]{article}
\usepackage[backend=biber, style=ieee]{biblatex}
\usepackage{parskip}
\addbibresource{references.bib}
\begin{document}
\author{Darcy Stewart N11977965}
\title{EGB322 Diesel Engine Report}
\maketitle
\pagebreak
\tableofcontents
\pagebreak
\section{Introduction}
A gas power cycle refers to a thermodynamic process where the energy absorbed as heat is converted to mechanical work \cite{Dincer}. The equipment that facilitate the gas power cycle are referred to as heat engines. Heat engines can be defined as a cyclical system that allows for the conversion of heat energy into mechanical work through the consumption of a fuel \cite{Basu}. The most popular example of the heat engine is the Internal Combustion Engine.
\par
The Internal Combustion Engine (ICE) have various defined types based around the style of cycle and fuel that is consumed. The stroke count of the engine refers to the number of strokes required to complete a single cycle. For a four-stroke engine the stroke purposes are as follows; intake, compression, expansion and exhaust \cite{Proctor}. While the stroke count has a high level of significance in determining the function of the ICE the fuel type also has a major impact on the function of the engine.
\par
Differing fuel types have specific methods of ignition suited to them. In the case of engines that use diesel as their fuel source the ignition method implemented is to compress the air in the cylinder of the piston to the point that the temperature of the air is beyond the ignition temperature of the diesel. Following this the fuel is injected and combustion occurs due to the air being above the ignition temperature, consequently expanding the gas and driving the piston \cite{Sun}. This differs from the likes of petrol engines that employ a spark based ignition system.
\par
Due to the differences in engines it is often difficult to compare the engines to one another as each cycle has different characteristics. To overcome this hurdle ideal cycles were developed. The purpose of an ideal cycle is to provide a simple processes that is theoretically completely reversible. To achieve this numerous assumptions must be made, consequently making it impossible for the actual measured cycles to match the idealised cycles due to irreversibilities and inefficiencies testing processes. They do however allow for simpler method to study the effects of certain parameters that affect the whole model. The idealised Air Standard Diesel Cycle has 4 processes; isentropic compression, isobaric heat addition, isentropic expansion and constant volume heat rejection that follow that order\cite{Proctor}. This report will employ the use of the idealised model and other key metrics to analyse the performance of a Four-Stroke Diesel engine compared to the theoretical idealised model.



\end{document}